\documentclass{article}

% packages
\usepackage{amsmath, amsthm, thmtools, amsfonts, amssymb, luacode, catchfile, tikzducks, hyperref, ifthen}
\ifcsname c@kobocompile\endcsname
	\usepackage[a5paper, total={1072pt, 1448pt}, margin=10pt, includeheadfoot]{geometry} % set page margins
\else
	\usepackage[a4paper, margin=50pt, includeheadfoot]{geometry}
\fi
\usepackage[shortlabels]{enumitem}
\usepackage[skip=3pt, indent=0pt]{parskip}

% language
\usepackage[bidi=basic, layout=tabular, provide=*]{babel}
\ifcsname c@english\endcsname
	\babelprovide[main, import]{english}
\else
	\babelprovide[main, import]{hebrew}
	\babelprovide{rl}
\fi
%\babelfont{rm}{Libertinus Serif}
\babelfont{rm}[Renderer=Harfbuzz]{Libertinus Serif}
\babelfont{sf}{Libertinus Sans}
\babelfont{tt}{Libertinus Mono}

% style
\AddToHook{cmd/section/before}{\clearpage}	% Add line break before section
\linespread{1.3}
\setcounter{secnumdepth}{0}		% Remove default number tags from sections, this won't do well with theorems
\AtBeginDocument{\setlength{\belowdisplayskip}{3pt}}
\AtBeginDocument{\setlength{\abovedisplayskip}{3pt}}
\graphicspath{ {../images/} }

% operators
\DeclareMathOperator\cis{cis}
\DeclareMathOperator\Sp{Sp}
\DeclareMathOperator\tr{tr}
\DeclareMathOperator\im{Im}
\DeclareMathOperator\re{Re}
\DeclareMathOperator\diag{diag}
\DeclareMathOperator*\lowlim{\underline{lim}}
\DeclareMathOperator*\uplim{\overline{lim}}
\DeclareMathOperator\rng{rng}
\DeclareMathOperator\Sym{Sym}
\DeclareMathOperator\Arg{Arg}
\DeclareMathOperator\Log{Log}
\DeclareMathOperator\dom{dom}
\DeclareMathOperator\supp{Supp}
\DeclareMathOperator\var{Var}
\DeclareMathOperator\cov{Cov}

% commands
%\renewcommand\qedsymbol{\textbf{מש''ל}}
%\renewcommand\qedsymbol{\fbox{\emoji{lizard}}}
\newcommand{\Aa}[0]{\mathcal{A}}
\newcommand{\Bb}[0]{\mathcal{B}}
\newcommand{\CC}[0]{\mathbb{C}}
\newcommand{\Cc}[0]{\mathcal{C}}
\newcommand{\EE}[0]{\mathbb{E}}
\newcommand{\FF}[0]{\mathbb{F}}
\newcommand{\Ff}[0]{\mathcal{F}}
\newcommand{\Ii}[0]{\mathcal{I}}
\newcommand{\Gg}[0]{\mathcal{G}}
\newcommand{\Ll}[0]{\mathcal{L}}
\newcommand{\Mm}[0]{\mathcal{M}}
\newcommand{\NN}[0]{\mathbb{N}}
\newcommand{\Nn}[0]{\mathcal{N}}
\newcommand{\PP}[0]{\mathbb{P}}
\newcommand{\Pp}[0]{\mathcal{P}}
\newcommand{\QQ}[0]{\mathbb{Q}}
\newcommand{\RR}[0]{\mathbb{R}}
\newcommand{\Rr}[0]{\mathcal{R}}
\newcommand{\Ss}[0]{\mathcal{S}}
\newcommand{\TT}[0]{\mathbb{T}}
\newcommand{\Uu}[0]{\mathcal{U}}
\newcommand{\Vv}[0]{\mathcal{V}}
\newcommand{\Ww}[0]{\mathcal{W}}
\newcommand{\ZZ}[0]{\mathbb{Z}}
\newcommand{\acts}[0]{\circlearrowright}
\newcommand{\explain}[2] {
	\begin{flalign*}
		 && \text{#2} && \text{#1}
	\end{flalign*}
}
\newcommand{\maketitleprint}[0]{ \begin{center}
	%\begin{tikzpicture}[scale=3]
	%	\duck[graduate=gray!20!black, tassel=red!70!black]
	%\end{tikzpicture}	
	\includegraphics[width=6cm]{cover}
\end{center}
}

% theorem commands
\newtheoremstyle{c_remark}
	{}	% Space above
	{}	% Space below
	{}% Body font
	{}	% Indent amount
	{\bfseries}	% Theorem head font
	{}	% Punctuation after theorem head
	{.5em}	% Space after theorem head
	{\thmname{#1}\thmnumber{ #2}\thmnote{ \normalfont{\text{(#3)}}}}	% head content
\newtheoremstyle{c_definition}
	{3pt}	% Space above
	{3pt}	% Space below
	{}% Body font
	{}	% Indent amount
	{\bfseries}	% Theorem head font
	{}	% Punctuation after theorem head
	{.5em}	% Space after theorem head
	{\thmname{#1}\thmnumber{ #2}\thmnote{ \normalfont{\text{(#3)}}}}	% head content
\newtheoremstyle{c_plain}
	{3pt}	% Space above
	{3pt}	% Space below
	{\itshape}% Body font
	{}	% Indent amount
	{\bfseries}	% Theorem head font
	{}	% Punctuation after theorem head
	{.5em}	% Space after theorem head
	{\thmname{#1}\thmnumber{ #2}\thmnote{ \text{(#3)}}}	% head content

\ifcsname c@english\endcsname
	\theoremstyle{plain}
	\newtheorem{theorem}{Theorem}[section]
	\newtheorem{lemma}[theorem]{Lemma}
	\newtheorem{proposition}[theorem]{Proposition}
	\newtheorem*{proposition*}{Proposition}
	%\newtheorem{corollary}[theorem]{אין חלופה עברית}

	\theoremstyle{definition}
	\newtheorem{definition}[theorem]{Definition}
	\newtheorem*{definition*}{Definition}
	\newtheorem{example}{Example}[section]
	\newtheorem{exercise}{Exercise}[section]

	\theoremstyle{remark}
	\newtheorem*{remark}{Remark}
	\newtheorem*{solution}{Solution}
	\newtheorem{conclusion}[theorem]{Conclusion}
	\newtheorem{notation}[theorem]{Notation}
\else
	\theoremstyle{c_plain}
	\newtheorem{theorem}{משפט}[section]
	\newtheorem{lemma}[theorem]{למה}
	\newtheorem{proposition}[theorem]{טענה}
	\newtheorem*{proposition*}{טענה}
	%\newtheorem{corollary}[theorem]{אין חלופה עברית}

	\theoremstyle{c_definition}
	\newtheorem{definition}[theorem]{הגדרה}
	\newtheorem*{definition*}{הגדרה}
	\newtheorem{example}{דוגמה}[section]
	\newtheorem{exercise}{תרגיל}[section]

	\theoremstyle{c_remark}
	\newtheorem*{remark}{הערה}
	\newtheorem*{solution}{פתרון}
	\newtheorem{conclusion}[theorem]{מסקנה}
	\newtheorem{notation}[theorem]{סימון}
\fi

% Questions related commands
\newcounter{question}
\setcounter{question}{1}
\newcounter{sub_question}
\setcounter{sub_question}{1}

\ifcsname c@english\endcsname
	\newcommand{\question}[1][0]{
		\ifthenelse{#1 = 0}{}{\setcounter{question}{#1}}
		\section{Question \arabic{question}}
		\addtocounter{question}{1}
		\setcounter{sub_question}{1}
	}

	\newcommand{\subquestion}[1][0]{
		\ifthenelse{#1 = 0}{}{\setcounter{sub_question}{#1}}
		\subsection{Part \alph{sub_question}}
		\addtocounter{sub_question}{1}
	}
\else
	\newcommand{\question}[1][0]{
		\ifthenelse{#1 = 0}{}{\setcounter{question}{#1}}
		\section{שאלה \arabic{question}}
		\addtocounter{question}{1}
		\setcounter{sub_question}{1}
	}

	\newcommand{\subquestion}[1][0]{
		\ifthenelse{#1 = 0}{}{\setcounter{sub_question}{#1}}
		\subsection{סעיף \localecounter{letters.gershayim}{sub_question}}
		\addtocounter{sub_question}{1}
	}
\fi

% import lua and start of document
\directlua{common = require ('../common')}

\GetEnv{AUTHOR}

% headers
\author{\AUTHOR}
\date\today

\title{פתרון מטלה 09 --- מבנים אלגבריים (2), 80446}

\begin{document}
\maketitle
\maketitleprint[red]

\question{}
תהי $L / K$ הרחבה סופית.
לכל $\alpha \in L$ נגדיר את הפונקציה $M_{\alpha} : L \to L$ על־ידי,
\[
	M_{\alpha}(x) = \alpha \cdot x
\]
בתרגול הגדרנו את הפונקציה $\tr_{L / K} : L \to K$ על־ידי,
\[
	\tr_{L / K}(\alpha)
	= \tr(M_{\alpha})
\]
ובאופן דומה את $N_{L / K} : L \to K$ על־ידי $N_{L / K}(\alpha) = \det(M_{\alpha})$.

יהי $\alpha \in L$ עם הפולינום המינימלי $f(x) = x^d + c_1 x^{d - 1} + \cdots + c_d$. \\
נראה בסעיפים הבאים כי,
\[
	\tr_{L / K}(\alpha)
	= \frac{[L : K]}{d},
	\qquad
	N_{L / K}(\alpha)
	= c_d^{[L : K] / d}
\]

\subquestion{}
נגדיר את הבסיס,
\[
	C
	= (1, x, \ldots, x^{d - 1})
\]
בסיס ל־$K(\alpha) / K$.

אם $B = (b_1, \ldots, b_t)$ בסיס ל־$L / K(\alpha)$ אז גם,
\[
	D
	= {( x^i \cdot b_j )}_{0 \le i < d, 1 \le j \le t}
\]
בסיס ל־$L / K$.

\subquestion{}
תהי $T : K(\alpha) \to K(\alpha)$ העתקה לינארית של כפל ב־$\alpha$. \\
נראה שמתקיים,
\[
	{[T]}_C
	= \begin{pmatrix}
		0 & \cdots & 0 & -c_d \\
		1 & 0 & \cdots & -c_{d - 1} \\
		  & \ddots & \ddots & \vdots \\
		0 & \cdots & 1 & -c_1
	\end{pmatrix}
\]
\begin{proof}
	נבחין כי $\alpha = x$ כווקטור מעל המרחב $K(\alpha)$.
	עבור $n < d$ מתקיים,
	\[
		T(c_n)
		= T x^n
		= \alpha x^n
		= x \cdot x^n
		= x^{n + 1}
	\]
	ונבחין במקרה המיוחד $n = d - 1$,
	\[
		x^d
		= -(c_1 x^{d - 1} + \cdots + c_d)
	\]
	ישירות מהגדרת השדה, נכתוב את המטריצה על־פי ערכי הווקטורים שמצאנו ונקבל,
	\[
		{[T]}_C
		= \begin{pmatrix}
			0 & \cdots & 0 & -c_d \\
			1 & 0 & \cdots & -c_{d - 1} \\
			  & \ddots & \ddots & \vdots \\
			0 & \cdots & 1 & -c_1
		\end{pmatrix} 
	\]
	בדיוק כפי שרצינו.
\end{proof}

\subquestion{}
נראה שעבור הבסיס $D$ מתקיים,
\[
	{[ M_{\alpha} ]}_D
	= \diag({[T]}_C, \ldots, {[T]}_C)
\]
ונסיק את הטענה.
\begin{proof}
	יהי $0 \le i < d$ ו־$1 \le j \le t$, אז,
	\[
		M_{\alpha} x^i b_j
		= \alpha \cdot x^i \cdot b_j
		= x^{i + 1} b_j
	\]
	ויחד עם הסעיף הקודם נקבל ישירות את המסקנה המבוקשת.

	עתה נוכל לחשב ישירות ונקבל,
	\[
		\tr_{L / K}(\alpha)
		= [L : K(\alpha)] \cdot (-c_1)
		= - \frac{[L : K]}{d} c_1
	\]
	ובאופן דומה מחוקי דטרמיננטות,
	\[
		N_{L / K}(\alpha)
		= {(\det {[T]}_C)}^{[L : K] / d}
		= -c_d^{[L : K] / d}
	\]
\end{proof}

\subquestion{}
נסיק שאם $\alpha_1, \ldots, \alpha_d$ הם הצמודים של $\alpha$ ב־$\overline{K}$ אז,
\[
	\tr_{L / K}(\alpha) = \frac{[L : K]}{d} \sum_{i = 1}^d \alpha_i
\]
וכן,
\[
	N_{N / K}(\alpha)
	= {\left(\prod_{i = 1}^d \alpha_i\right)}^{[L : K] / d}
\]
\begin{proof}
	מתקיים,
	\[
		f(x) = \prod_{i = 1}^d (x - \alpha_i)
	\]
	ולכן,
	\[
		c_1
		= \sum_{i = 1}^d \alpha_i
	\]
	מחישוב שנעשה במטלות קודמות, נסיק את החלק של הטענה ישירות ממסקנת הסעיף הקודם.

	מאותה שאלה ממטלות קודמות גם,
	\[
		c_d
		= \prod_{i = 1}^d \alpha_i
	\]
	והחלק השני של הטענה נובע.
\end{proof}

\question{}
יהי $F$ שדה ונניח ש־$L = F(t_1, \ldots, t_n)$ עם הפעולה של $S_n$ המוגדרת על־ידי,
\[
	\sigma \cdot P
	= P(t_{\sigma(1)}, \ldots, t_{\sigma(n)})
\]
נסמן ב־$s_1, \ldots, s_n$ את הפולינומים הסימטריים האלמנטריים ב־$t_1, \ldots, t_n$,
כלומר הפולינומים המקיימים,
\[
	\prod_{i = 1}^n (x - t_i)
	= x^n - s_1 x^{n - 1} + \cdots + {(-1)}^n s_n
\]
בכל סעיף נגדיר איבר $P \in L^{S_n}$ ונבטא אותו באמצעות הפולינומים הסימטריים האלמנטריים.

\subquestion{}
נגדיר $P = t_1^3 + \cdots + t_n^3$.
\begin{solution}
	נסמן,
	\[
		Q = \sum_{\substack{1 \le i, j \le n \\ i \ne j}} t_i^2 t_j
	\]
	מחישוב ישיר מתקיים,
	\[
		s_1^3
		= P + s_3 + Q
		\iff P = s_1^3 - s_3 - Q
	\]
	ונבחין כי גם מתקיים,
	\[
		s_1 s_2
		= s_3 + Q
	\]
	ולכן,
	\[
		Q = s_1 s_2 - s_3
	\]
	ובהתאם גם,
	\[
		P
		= s_1^3 - s_1 s_2
	\]
\end{solution}

\subquestion{}
נגדיר את הפולינום,
\[
	P
	= \sum_{\substack{1 \le i, j \le n \\ i \ne j}} t_i^2 t_j
\]
\begin{solution}
	בהתאם לסעיף הקודם,
	\[
		P = s_1 s_2 - s_3
	\]
\end{solution}

\subquestion{}
עבור $n = 3$ נגדיר,
\[
	P = {((t_1 - t_2) (t_1 - t_3) (t_2 - t_3))}^2
\]
\begin{solution}
	\[
		(t_1 - t_2) (t_1 - t_3) (t_2 - t_3)
		= 4 s_3 + (t_1^2 t_2 - t_1^2 t_3 - t_2^2 t_1 + t_2^2 t_3 + t_3^2 t_1 - t_3^2 t_2)
		= s_2^2 - 2 s_3 s_1
	\]
	מבדיקה ישירה, ולכן גם,
	\[
		P = {(s_2^2 - s s_3 s_1)}^2
	\]
\end{solution}

\question{}
\subquestion{}
תהי $L / K$ הרחבת גלואה כך ש־$G = \gal(L / K)$.
נראה שלכל $H \le G$ מתקיים $[G : H] = [L^H : K]$.
\begin{proof}
	אנו כבר יודעים כי $[L : K] = |G|$ וכן שמתקיים,
	\[
		[G : H]
		= \frac{|G|}{|H|}
		= \frac{[L : K]}{|H|}
		= \frac{[L : L^H] \cdot [L^H : K]}{|H|}
	\]
	ולכן ננסה לחשב את $[L : L^H]$. \\
	בהרצאה ראינו את הטענה כי $H = \gal(L / L^H)$ ולכן ממשפט אפיון גודל חבורת האוטומורפיזמים, $[L : L^H] = |H|$
	נקבל אם כך,
	\[
		[G : H]
		= \frac{[L : L^H] \cdot [L^H : K]}{|H|}
		= \frac{[L : L^H] \cdot |H|}{|H|}
		= [L^H : K]
	\]
	בדיוק כפי שרצינו.
\end{proof}

\subquestion{}
נוכיח של־$A_4$ אין תת־חבורה מאינדקס 2,
ונסיק שאם $A_4 \simeq \gal(L / K)$ אז אין הרחבת ביניים $L / F / K$ כך ש־$[F : K] = 2$.
\begin{proof}
	אנו יודעים ש־$|S_4| = 4! = 24$ וכן ש־$|A_4| = 4! / 2 = 12$.
	אילו הייתה $H \le A_4$ תת־חבורה מאינדקס 2, אז $|H| = 6$ בדיוק.
	בהתאם אנו גם יודעים כי $H \simeq \ZZ / 6\ZZ$ או $H \simeq S_3$ בלבד.
	אם $H \simeq \ZZ / 6\ZZ$ אז יש תמורה $\sigma \in A_4$ מסדר $6$,
	אבל כל תמורה כזו היא מסדר 3 כמחזור מגודל 3, או שני מחזורים מסדר 2, ובהתאם סדר $\sigma$ הוא $2$ או $3$ בלבד, ולכן $H \simeq S_3$.
	בהתאם חייב להיות $\sigma \in S_3$ מסדר $3$ ו־$\tau \in S_3$ מסדר $2$.
	בהתאם $\{ \sigma, \sigma^2 \} \subseteq S_3$ וכן $S_3 \ni \tau^{-1} \sigma \tau \notin \{ \sigma, \sigma^2, e \}$ ונקבל ישירות ש־$|H| > 6$ בסתירה.
	לכן לא קיים $H$ כזה.

	נניח ש־$L / K$ הרחבת גלואה כך ש־$A_4 \simeq \gal(L / K)$.
	נניח בשלילה ש־$L / F / K$ כך ש־$[F : K] = 2 = |H|$ עבור $H \le A_4$ בהתאם לסעיף הקודם.
	אבל ראינו עתה כי לא קיים $H$ כזה, ולכן בהכרח גם לא קיים $F$ כזה.
\end{proof}

\end{document}
